\documentclass[12pt,a4paper]{article}

\usepackage[utf8]{inputenc}
\usepackage[T1]{fontenc}
\usepackage[a4paper,margin=3cm]{geometry}
\usepackage{amsmath}
\usepackage{booktabs}
\usepackage{caption}
\usepackage{graphicx}
\usepackage{float}
\usepackage{xcolor}
\usepackage{listings}
\usepackage{hyperref}

% Gunakan font reguler yang tersedia pada distribusi LaTeX minimal.
% Bobot tebal dipetakan ke bobot reguler agar Tectonic minimal tetap bisa
% menghasilkan PDF tanpa font bitmap tambahan.
\renewcommand{\bfdefault}{m}

\hypersetup{
    colorlinks=true,
    linkcolor=blue,
    urlcolor=blue,
    citecolor=blue
}

\graphicspath{{screenshots/}}

% Menampilkan gambar jika tersedia. Jika belum ada, tampilkan kotak
% penanda agar dokumen tetap dapat dikompilasi.
\newcommand{\safeincludegraphics}[2][]{%
    \IfFileExists{screenshots/#2}{%
        \includegraphics[#1]{#2}%
    }{%
        \fbox{%
            \begin{minipage}{0.9\textwidth}
                \centering
                \vspace{1.5cm}
                Gambar belum tersedia: \texttt{#2}
                \vspace{1.5cm}
            \end{minipage}%
        }%
    }%
}

\lstdefinestyle{pythonstyle}{
    language=Python,
    basicstyle=\ttfamily\small,
    keywordstyle=\color{blue},
    commentstyle=\color{gray},
    stringstyle=\color{orange!80!black},
    numbers=left,
    numberstyle=\tiny\color{gray},
    frame=single,
    breaklines=true,
    showstringspaces=false
}

\title{\textbf{Peningkatan Kualitas Citra dengan Histogram Equalization}}
\author{
    Nama: Juyus Muhammad Adinulhaq\\
    NIM: 572923\\
    Mata Kuliah: Pengolahan dan Analisis Citra Digital (MII226205)
}
\date{\today}

\begin{document}

\maketitle
\tableofcontents
\newpage

\section{Pendahuluan}

\subsection{Latar Belakang}
Citra digital merupakan salah satu bentuk informasi visual yang banyak digunakan dalam berbagai bidang, seperti dokumentasi, pengolahan medis, pengawasan, dan analisis objek. Kualitas citra sangat dipengaruhi oleh tingkat kontrasnya. Citra dengan kontras rendah memiliki rentang intensitas piksel yang sempit sehingga perbedaan antara area terang dan area gelap menjadi kurang terlihat.

Kontras rendah dapat disebabkan oleh kondisi pencahayaan yang kurang baik, bayangan, kabut, keterbatasan sensor kamera, atau proses pengambilan citra yang tidak optimal. Akibatnya, detail pada citra menjadi sulit diamati dan informasi visual yang terkandung di dalamnya tidak dapat terlihat secara maksimal.

Salah satu metode yang dapat digunakan untuk meningkatkan kontras citra adalah Histogram Equalization (HE). Metode ini memanfaatkan distribusi intensitas piksel pada histogram citra untuk memetakan nilai intensitas lama ke nilai intensitas baru. Dengan demikian, rentang intensitas citra dapat diperluas sehingga detail pada area gelap maupun terang menjadi lebih terlihat.

Pada laporan ini, Histogram Equalization diterapkan secara manual pada citra grayscale. Selain itu, hasil implementasi manual dibandingkan dengan hasil dari fungsi \texttt{cv2.equalizeHist()} untuk memastikan bahwa proses yang dilakukan telah sesuai.

\subsection{Tujuan}
Tujuan dari eksperimen ini adalah sebagai berikut:

\begin{enumerate}
    \item Menerapkan metode Histogram Equalization untuk meningkatkan kontras citra.
    \item Memahami tahapan Histogram Equalization, mulai dari perhitungan histogram, CDF, normalisasi, hingga pemetaan nilai piksel.
    \item Mengimplementasikan Histogram Equalization secara manual tanpa hanya menggunakan fungsi bawaan OpenCV.
    \item Membandingkan hasil implementasi manual dengan hasil dari fungsi \texttt{cv2.equalizeHist()}.
    \item Menganalisis perubahan citra berdasarkan tampilan visual, histogram, kurva CDF, dan statistik citra.
\end{enumerate}

\subsection{Citra yang Digunakan}
Citra yang digunakan dalam eksperimen ini adalah citra bernama \texttt{[nama\_file\_citra]} yang diperoleh dari \texttt{[sumber citra, misalnya Google atau dataset tertentu]}. Citra tersebut memiliki format \texttt{[PNG/JPG]} dengan ukuran sebesar \texttt{[lebar] $\times$ [tinggi]} piksel.

Sebelum diproses menggunakan Histogram Equalization, citra dikonversi terlebih dahulu menjadi citra grayscale. Konversi ini dilakukan agar setiap piksel hanya memiliki satu nilai intensitas dengan rentang 0 sampai 255. Nilai 0 menunjukkan warna hitam, sedangkan nilai 255 menunjukkan warna putih.

Citra ini dipilih karena memiliki tingkat kontras yang relatif rendah. Berdasarkan histogram citra, sebagian besar nilai intensitas piksel menumpuk pada rentang \texttt{[contoh: 60--150]}. Penumpukan nilai pada rentang intensitas yang sempit menunjukkan bahwa citra belum memanfaatkan seluruh rentang intensitas 0 sampai 255 secara optimal. Akibatnya, beberapa detail pada area gelap atau terang terlihat kurang jelas.

Melalui Histogram Equalization, nilai-nilai intensitas yang menumpuk tersebut dipetakan ke rentang yang lebih luas. Proses ini diharapkan dapat meningkatkan perbedaan antara area gelap dan terang sehingga detail citra menjadi lebih mudah diamati.

\section{Dasar Teori}

\subsection{Histogram Citra}
Jelaskan pengertian histogram citra dan hubungan distribusi intensitas dengan kontras.

\subsection{Histogram Equalization}
Histogram Equalization mengubah nilai intensitas berdasarkan histogram kumulatif atau \textit{cumulative distribution function} (CDF). Tahapan umumnya adalah menghitung histogram $h(i)$, menghitung CDF, melakukan normalisasi, kemudian memetakan nilai piksel.

Rumus pemetaan nilai intensitas dituliskan sebagai berikut:

\begin{equation}
    \text{baru}(i) = \operatorname{round}\left(
    \frac{\operatorname{cdf}(i)-\operatorname{cdf}_{\min}}
    {N-\operatorname{cdf}_{\min}} \times 255
    \right)
\end{equation}

Keterangan:
\begin{itemize}
    \item $i$ adalah nilai intensitas piksel.
    \item $\operatorname{cdf}(i)$ adalah nilai CDF pada intensitas $i$.
    \item $\operatorname{cdf}_{\min}$ adalah nilai CDF minimum yang bukan nol.
    \item $N$ adalah jumlah piksel pada citra.
\end{itemize}

Jelaskan dengan bahasa sendiri mengapa pemetaan tersebut dapat meningkatkan kontras. Hubungkan penjelasan dengan nilai intensitas yang sering muncul dan kecuraman CDF.

\section{Metodologi dan Implementasi}

\subsection{Alur Kerja}
Alur eksperimen yang dilakukan adalah:

\begin{enumerate}
    \item Membaca citra asli.
    \item Mengubah citra menjadi grayscale.
    \item Menghitung histogram.
    \item Menghitung CDF.
    \item Menormalisasi CDF.
    \item Memetakan setiap piksel ke nilai intensitas baru.
    \item Membandingkan hasil manual dengan implementasi OpenCV.
\end{enumerate}

\subsection{Implementasi Manual}
Masukkan kode implementasi manual Histogram Equalization berikut atau lampirkan kode lengkap pada bagian lampiran.

\begin{lstlisting}[style=pythonstyle,caption={Implementasi manual Histogram Equalization}]
# Tempel kode implementasi manual HE dari Google Colab di sini.
# Kode harus menghitung histogram, CDF, normalisasi,
# dan pemetaan piksel tanpa hanya memanggil cv2.equalizeHist().
\end{lstlisting}

\begin{figure}[H]
    \centering
    \safeincludegraphics[width=0.95\textwidth]{kode-he-manual.png}
    \caption{Kode implementasi Histogram Equalization secara manual.}
\end{figure}

\subsection{Verifikasi terhadap OpenCV}
Jelaskan metode verifikasi dan nilai selisih antara hasil implementasi manual dengan hasil dari \texttt{cv2.equalizeHist()}.

\begin{figure}[H]
    \centering
    \safeincludegraphics[width=0.85\textwidth]{verifikasi-opencv.png}
    \caption{Verifikasi hasil manual terhadap \texttt{cv2.equalizeHist()}.}
\end{figure}

\section{Hasil dan Pembahasan}

\subsection{Perbandingan Visual}
Jelaskan perubahan visual citra asli setelah Histogram Equalization, terutama perubahan kontras dan detail yang terlihat.

\begin{figure}[H]
    \centering
    \safeincludegraphics[width=0.95\textwidth]{perbandingan-visual.png}
    \caption{Perbandingan citra asli dan citra hasil Histogram Equalization.}
\end{figure}

\subsection{Perbandingan Histogram}
Jelaskan perubahan bentuk histogram sebelum dan sesudah HE. Hubungkan perubahan tersebut dengan penyebaran nilai intensitas.

\begin{figure}[H]
    \centering
    \safeincludegraphics[width=0.9\textwidth]{histogram-overlay.png}
    \caption{Histogram citra sebelum dan sesudah Histogram Equalization.}
\end{figure}

\subsection{Kurva CDF}
Jelaskan hubungan antara kecuraman kurva CDF dan seberapa jauh nilai piksel direnggangkan oleh proses pemetaan.

\begin{figure}[H]
    \centering
    \safeincludegraphics[width=0.9\textwidth]{kurva-cdf.png}
    \caption{Kurva CDF citra sebelum dan sesudah Histogram Equalization.}
\end{figure}

\subsection{Statistik Kuantitatif}
Masukkan nilai mean, standar deviasi, dan entropy sebelum serta sesudah HE. Interpretasikan perubahan setiap nilai.

\begin{table}[H]
    \centering
    \caption{Perbandingan statistik citra sebelum dan sesudah HE.}
    \label{tab:statistik}
    \begin{tabular}{lcc}
        \toprule
        Statistik & Sebelum HE & Sesudah HE \\
        \midrule
        Mean & [isi] & [isi] \\
        Standar deviasi & [isi] & [isi] \\
        Entropy & [isi] & [isi] \\
        \bottomrule
    \end{tabular}
\end{table}

\begin{figure}[H]
    \centering
    \safeincludegraphics[width=0.85\textwidth]{statistik.png}
    \caption{Output statistik mean, standar deviasi, dan entropy.}
\end{figure}

\subsection{Uji Tambahan}

\subsubsection{HE pada Citra dengan Kontras Tinggi}
Jelaskan hasil penerapan HE pada citra yang sejak awal memiliki kontras tinggi. Bahas apakah citra berubah sedikit atau muncul artefak, lalu hubungkan dengan kurva CDF.

\begin{figure}[H]
    \centering
    \safeincludegraphics[width=0.95\textwidth]{uji-kontras-tinggi.png}
    \caption{Hasil uji HE pada citra dengan kontras tinggi.}
\end{figure}

\subsubsection{HE pada Citra Berwarna melalui HSV}
Jelaskan perbedaan penerapan HE pada channel V dalam ruang warna HSV dibandingkan penerapan HE pada channel R, G, dan B secara terpisah.

\begin{figure}[H]
    \centering
    \safeincludegraphics[width=0.95\textwidth]{uji-warna-hsv.png}
    \caption{Perbandingan penerapan HE pada citra berwarna melalui channel V HSV.}
\end{figure}

\section{Kesimpulan}
Ringkas kapan Histogram Equalization efektif dan kapan kurang efektif berdasarkan hasil eksperimen. Sebutkan juga keterbatasannya, terutama sifatnya yang global dan tidak menyesuaikan diri terhadap area lokal citra.

\section{Referensi}
Tuliskan semua sumber teori, dokumentasi, artikel, atau buku yang digunakan.

\begin{thebibliography}{9}
    \bibitem{opencv}
    OpenCV Documentation, ``Histogram Equalization,''
    \url{https://docs.opencv.org/}.

    \bibitem{referensi-tambahan}
    [Tambahkan referensi lain yang digunakan].
\end{thebibliography}

\appendix
\section{Kode Lengkap}
Lampirkan kode lengkap dari notebook Google Colab atau tautan menuju notebook.

\begin{lstlisting}[style=pythonstyle,caption={Kode lengkap eksperimen}]
# Tempel kode lengkap di sini atau tuliskan tautan notebook Colab.
\end{lstlisting}

\end{document}
